% Gomez--Venegas: supplemental verification note (22 November 2026).
% Scope and caveats are deliberate; do not cite this as a classification theorem.
\section{Boundary of the $\mathbb D$--$\mathbb G$ verification}
The proof of the $h=1$ ozone theorem for $\mathbb G$ in
\texttt{DG\_THEOREM\_NOTE.tex} depends on the full-center theorem
\texttt{thm:G} in the working manuscript. A change of twist $h$ may change
both the center and the ozone group. Thus the full $\mathbb G(p,f,h)$
classification is \emph{not} certified. For $\mathbb D(p,1)$ the necessary
central constraints leave 32 diagonal candidates at $p=i$, whereas the
historical table reports ozone order 16. Until additional central
constraints are established, the $\mathbb D$ classification is also open.

\section{An exact quotient-conic certificate for $\mathbb Q$ and $\mathbb V$}
Let $n$ be the distinguished normal generator and $t=n^{-1}x'$ as in
Table~\texttt{tab:Mt} of the working manuscript. In each of the families
$\mathbb Q$ and $\mathbb V$, writing $Y=(y_1,y_2)^T$, the relations take
$Yt=M(t)Y$, where
\[
M_{\mathbb Q}(t)=\begin{pmatrix}-1&1-t\\1+t&1\end{pmatrix},\qquad
M_{\mathbb V}(t)=\begin{pmatrix}1&t\\1-t&0\end{pmatrix}.
\]
The characteristic polynomials are
\[
\chi_{\mathbb Q}(t,\mu)=\mu^2+t^2-2,\qquad
\chi_{\mathbb V}(t,\mu)=\mu^2-\mu+t^2-t.
\]
The associated Vieta transformations are
\[
\phi_{\mathbb Q}(t,\mu)=(\mu,-t),\qquad
\phi_{\mathbb V}(t,\mu)=(\mu,1-t).
\]
\begin{proposition}
Each of these transformations has exact order four and preserves the
respective affine conic. Both are rotations by a quarter-turn in
suitable centered coordinates: for $\mathbb Q$ about $(0,0)$, for
$\mathbb V$ about $(\frac12,\frac12)$.
\end{proposition}
\begin{proof}
The Jacobian matrix of each transformation is
$J=\left(\begin{smallmatrix}0&1\\-1&0\end{smallmatrix}\right)$,
so $J^2=-I$ and $J^4=I$. The stated centers are fixed points.
Substitution gives $\chi(t,\mu)=\chi(\phi(t,\mu))$; and
$\phi^2\ne\mathrm{id}$ on the coordinate ring of the conic since
that conic is positive-dimensional in characteristic zero. Hence both
maps have exact order four.
\end{proof}
This is a certificate of the \emph{dynamical reduction}, not yet a
certificate for the ozone group. Determining $\operatorname{Oz}$ still
requires a complete generating set of the affine center, the degree-one
matrix equations, and a proof of exhaustivity. Family $\mathbb R$
requires its separate degree-two-normal-element presentation; a
matrix $M(t)$ must not be inferred from the data for $\mathbb Q$ or
$\mathbb V$.
